\section{Introduction}\label{sec:introduction}

Recent evidence associates generative-AI exposure with weaker employment growth among workers in their early twenties.  Interpretation is difficult because exposure is constructed rather than observed, high- and low-exposure jobs differ broadly, and the period around ChatGPT's release includes several aggregate disruptions.

This paper asks which occupational comparisons support the public Current Population Survey (CPS) pattern.  I rebuild monthly employment stocks from January 2017 through July 2026 and classify occupations with the GPT-4 beta score of \citet{eloundou2024gpts}.  The target is the post-December-2022 change in the log conditional-mean employment-stock ratio of ages 22--25 to ages 26--65 in the highest relative to the lowest exposure quintile, conditional on intermediate quintiles and preperiod Webb software exposure.  I examine broad-family paths, direct within-family support, preperiod computer use, dynamics, and linked-CPS flows.  Exposure measures susceptibility rather than adoption; January 2023 is chronology, not an instrument.

The reconstructed coefficient is $-0.132$ log points, with 95-percent occupation wild-score interval $[-0.221,-0.044]$; it is neither an individual probability nor an employer hiring rate.  With two-digit occupation-family young-relative monthly paths, it is $-0.022$ ($[-0.161,0.117]$).  Their paired difference is $0.110$ ($[0.011,0.210]$).  This distinguishes conditional estimands without making a causal decomposition: family paths may absorb AI-related and non-AI-related evolution.

Direct support is thin.  Four families span both exposure tails; their 29 occupations represent 5.03 percent of preperiod stock, and the interval includes large effects of either sign.  The conditioned comparison otherwise relies on intermediate quintiles and common-coefficient restrictions, so its near-zero estimate is not economic equivalence or evidence of no AI effect.

Adding a pre-2022 O*NET computer-use measure makes the common-support coefficient more negative; one control therefore does not resolve computerization.  Unrestricted quarterly coefficients reject preperiod equality, while linked-flow intervals include zero.  The data neither isolate an AI counterfactual nor attribute the stock pattern to exit, switching, or entry allocation.

The paper shows that a sizeable public-CPS association is identified predominantly by comparisons spanning broad occupation families, quantifies the resolution of direct and flow evidence, and records calendar and taxonomy failures that silently alter support.  This localization tells future adoption or quasi-experimental studies which comparisons they must replace rather than inherit.  It narrows rather than overturns the motivating evidence: it neither replicates proprietary payroll outcomes nor identifies AI, computerization, or pandemic recovery as the cause.
